\documentclass{article}
\usepackage[utf8]{inputenc}
\usepackage{amsmath}
\usepackage{geometry}
\geometry{margin=2cm}
\begin{document}

\section*{Questão 139 — ENEM 2024 — Matemática}

Um artesão utiliza dois tipos de componentes, \(X\) e \(Y\), nos enfeites que produz. Ele sempre compra todos os componentes em uma mesma loja. A tabela a seguir apresenta os preços dos componentes nas lojas I e II:

\begin{center}
\begin{tabular}{|c|c|c|}
\hline
\rowcolor[HTML]{F4CCCC} 
\textbf{Lojas} & \textbf{X (R\$)} & \textbf{Y (R\$)} \\
\hline
I & 3,00 & 1,00 \\
\hline
II & 2,00 & 4,00 \\
\hline
\end{tabular}
\end{center}

Cada enfeite é formado por 2 unidades do componente \(X\) e 1 unidade do componente \(Y\). O artesão decidirá onde comprar os componentes para obter o menor custo.

\section*{Solução técnica}

Sabendo as quantidades, calculamos o custo total em cada loja:

\[
C_{I} = 2 \times 3 + 1 \times 1 = 6 + 1 = R\$ 7,00,
\]
\[
C_{II} = 2 \times 2 + 1 \times 4 = 4 + 4 = R\$ 8,00.
\]

Portanto, a compra na loja I é a mais econômica.

\section*{Explicação didática}

Para escolher a loja, calculamos o custo total dos componentes necessários para um enfeite e comparamos os valores. Lembre-se sempre de multiplicar o preço unitário pela quantidade para cada componente e somar para obter o custo total.

\section*{Análise dos distratores}

O erro comum nas alternativas incorretas envolve tomar o preço unitário isoladamente, esquecer de multiplicar pela quantidade ou calcular erroneamente o valor total.

\section*{Perfil da questão}

Questão de dificuldade média que exige habilidade em operações com números racionais, interpretação de tabelas e raciocínio quantitativo.

\end{document}